% Generated draft: requires mathematical and source review.
\documentclass[11pt]{article}
\usepackage[margin=25mm]{geometry}
\usepackage{fontspec}
\setmainfont{Latin Modern Roman}
\usepackage{amsmath,amssymb,mathtools,unicode-math}
\setmathfont{Latin Modern Math}
\usepackage{enumitem,needspace,xurl,hyperref}
\setlength{\parindent}{0pt}
\setlength{\parskip}{5pt}
\setlength{\emergencystretch}{3em}
\newcommand{\N}{\mathbb{N}}
\newcommand{\Z}{\mathbb{Z}}
\newcommand{\Q}{\mathbb{Q}}
\newcommand{\R}{\mathbb{R}}
\newcommand{\C}{\mathbb{C}}
\newcommand{\F}{\mathbb{F}}
\newcommand{\A}{\mathbb{A}}
\newcommand{\PP}{\mathbb{P}}
\newcommand{\E}{\mathbb{E}}
\newcommand{\eps}{\varepsilon}
\newcommand{\dd}{\,\mathrm{d}}
\begin{document}
\section*{UnsolvedMath ID: 30006010}
\textbf{Source problem code:} catalogue entry
\subsection{Definitions and mathematical statement}
The supplied record does not specify enough of the ambient framework to give a fully determined mathematical conjecture. The following states its requested properties while retaining those omissions explicitly.

Let \(G\) be a group in the intended, unspecified geometric Langlands setting, and let \(\operatorname{Bun}_G\) denote the corresponding moduli object of \(G\)-bundles. Write \(\Psi_A(G)\) for the collection of Arthur parameters, also called A-parameters, in that setting. The base field, underlying curve or geometric object, definition of an Arthur parameter, coefficient theory, and sheaf category are not supplied. Thus these symbols are placeholders for data that must be recovered from the source before the statement is mathematically complete.

The requested conjectural assertion is that for every \(\psi\in\Psi_A(G)\), there exists a sheaf object \(F_\psi\) on \(\operatorname{Bun}_G\) satisfying the following properties:
\begin{enumerate}
\item \(F_\psi\) has the predicted variant of the Hecke eigenproperty associated with \(\psi\). The supplied excerpt does not specify the Hecke functors, eigenvalue objects, shifts, or compatibility conditions, so no precise eigen-isomorphism can be stated from this record.
\item \(F_\psi\) is perverse for the intended perverse \(t\)-structure, which is not specified in the record.
\item If \(G\) is semisimple, then for every basic index \(b\), the restriction \(i_b^*F_\psi\) to the corresponding basic stratum has the predicted decomposition into the Arthur packet for the associated pure inner form \(G_b\). Here \(i_b\) denotes the inclusion of that stratum. The supplied source asserts, in its intended setting, that \(G_b\) is a pure inner form of \(G\) for every such \(b\).
\end{enumerate}
The damaged source display indexes the proposed decomposition by \(\pi\in\Pi_{\iota\circ\psi}(G_b(E))\). It does not define \(E\), \(\iota\), or this packet notation. Its remaining decomposition and pairing notation is corrupted, so the summands, multiplicities, gradings, shifts, and meaning of the asserted isomorphism cannot be recovered reliably. No exact decomposition formula is asserted here. In particular, the clean statement's phrase ``predicted decomposition'' remains a source-dependent requirement rather than a specified mathematical identity.
\subsection{Short English statement}
For every Arthur parameter, the problem asks for an associated perverse Hecke eigensheaf on \(\operatorname{Bun}_G\). When \(G\) is semisimple, its restrictions to basic strata should have the predicted Arthur-packet decompositions for the corresponding pure inner forms. The supplied record omits essential ambient definitions and contains a damaged decomposition formula; this is therefore an incomplete conjectural specification requiring source recovery and human mathematical review.
\subsection{Sources}
\begin{itemize}
\item Supplied local UnsolvedMath registry record, permanent ID \(30006010\), title: Hecke Eigensheaves for Arthur Parameters. Exact source problem code: \texttt{catalogue entry}. Collection: Oberwolfach Reports Open Problems. No individual proposer is identified; the supplied proposed year is 2024.
\item Supplied source citation: \emph{Arithmetic Geometry (2025)}, Oberwolfach Reports, DOI \texttt{10.4171/owr/2024/33}, pp.\ 1855--1912. \url{https://doi.org/10.4171/owr/2024/33}. The record describes these reports as documenting mathematical workshops and their open-problem discussions. The citation title gives 2025, whereas the DOI and proposed-year metadata give 2024; the record does not explain this discrepancy.
\item Supplied literature source, labeled Bertoloni Meli--Koshikawa--Leake: \url{https://arxiv.org/abs/2409.06006}. No title or supporting passage is supplied.
\item Supplied literature source, labeled Hansen--Mann: \url{https://arxiv.org/abs/2606.00983}. No title or supporting passage is supplied.
\item Literature assessment dated exactly \texttt{2026-08-22}: supplied dated curation, not independently verified. The record labels the problem partially solved and reports that candidate sheaves are constructed on the spectral side for arbitrary A-parameters, while the expected perverse, packet, and automorphic-side properties are not established in general. Its background explicitly calls this a dated literature triage derived from the supplied OWR curation and requests independent verification of time-sensitive status and cited claims.
\item The registry labels the statement \texttt{corrected\_verified}, with the qualification that only clear formula delimiters, notation, grammar, and typography were repaired and that mathematical content and scope were preserved. This supplied editorial label does not constitute independent mathematical verification of the formulation or its status.
\end{itemize}
\end{document}
